\section{SmallWood proof-stack details}
\label{app:smallwood}
\label{app:smallwood-details}

This appendix records the proof-stack geometry for the compiled relations of
Section~\ref{sec:smallwood-model}. It specifies committed witness-row families,
public lifts, PACS splits, and the PCS/LVCS/DECS opening structure used by the
prover and verifier. It does not strengthen the statement layer beyond
Section~\ref{sec:smallwood-model}; in particular, coefficient-domain
consequences are obtained exactly by the theorem suite stated there.

\newcommand{\NTT}{\mathrm{NTT}}

\subsection{Deferred row inventories and bridge formulas}
\label{subsec:smallwood-row-inventories}

Section~\ref{sec:smallwood-model} groups the witness rows semantically. We record here the exact compiled surfaces and
bridge formulas used by the issuance and showing relations.

\paragraph{Issuance witness surface.}
The committed issuance witness consists of carrier rows for the signed message lanes and the linear-side
randomness, together with decoded source rows
\[
M,\ K,\ S,\ E,\ Z.
\]
The exact carrier-to-source bridges decode the committed message/key lanes and the linear-side randomness lanes
into these source rows. The aligned commitment is then enforced directly by
\[
\vecc = \mA_m M + \mA_k K + \mA_s S + E,
\]
and the inverse witness is enforced by
\[
(B_3-R_1)\Had Z = 1.
\]
Finally, the public issuance target is tied to the same committed witness by the aligned identity
\[
T = B_0 + U\vecc + Z.
\]
Thus the pre-sign statement binds one coherent hidden witness to the commitment, the inverse witness, and the
public target.

\paragraph{Showing witness surface.}
The committed showing witness consists of:
\begin{enumerate}[leftmargin=1.25em]
    \item the signature limb rows \(L_{c,b,\ell}\);
    \item the carrier rows \(C^M\) and \(C^{\mathrm{ctr}}\);
    \item the source rows
    \[
    T,\ M,\ K,\ R_0,\ R_1,\ Z;
    \]
    \item the replay-image block families
    \[
    \widehat T_0,\ldots,\widehat T_{B_{\mathrm{rep}}-1},
    \qquad
    \widehat m_0,\ldots,\widehat m_{B_{\mathrm{rep}}-1},
    \]
    \[
    \widehat R_{0,0},\ldots,\widehat R_{0,B_{\mathrm{rep}}-1},
    \qquad
    \widehat R_{1,0},\ldots,\widehat R_{1,B_{\mathrm{rep}}-1};
    \]
    \item the PRF companion rows.
\end{enumerate}
The cert-side signature witness \(u\) is reconstructed from the limb surface by the public digit-collapse map of the
shortness gadget and is tied to the source row \(T\) by the source-side equation \(T=\mA u\). The carrier-to-source bridges are
\[
M=\mathrm{Dec}_M^{(1)}(C^M),
\qquad
K=\mathrm{Dec}_M^{(2)}(C^M),
\]
\[
R_0=\mathrm{Dec}_{\mathrm{ctr}}^{(1)}(C^{\mathrm{ctr}}),
\qquad
R_1=\mathrm{Dec}_{\mathrm{ctr}}^{(2)}(C^{\mathrm{ctr}}).
\]
The source-side inverse witness satisfies
\[
(B_3-R_1)\Had Z = 1.
\]
For every replay block \(b\in\{0,\ldots,B_{\mathrm{rep}}-1\}\), the exact source-to-replay bridges are
\[
\widehat T_b = \mathcal F_{\mathrm{rep},b}(T),
\qquad
\widehat m_b = \mathcal F_{\mathrm{rep},b}(M+K),
\]
\[
\widehat R_{0,b} = \mathcal F_{\mathrm{rep},b}(R_0),
\qquad
\widehat R_{1,b} = \mathcal F_{\mathrm{rep},b}(R_1).
\]

\subsection{Domains and packing}
Fix a base field \(\Fq\), an evaluation set \(\Omega\subset\Fq\) of size
\(s\), and a disjoint tail set \(\Omega'\subset\Fq\setminus\Omega\) of size
\(\ell\). A witness row is interpolated into a polynomial of degree
\(\le s+\ell-1\) whose values on \(\Omega\) are the packed coordinates and
whose values on \(\Omega'\) are random masks. The prover commits to these
explicit-domain row polynomials, and the verifier later replays the active
constraint residuals from openings of the same polynomials.

\subsection{DECS (Degree-Enforcing Commitment Scheme)}
DECS commits to evaluations of witness polynomials and enforces that they arise
from low-degree polynomials rather than arbitrary vectors. In our
instantiation, commitments are implemented by a Merkle tree over an NTT-sized
domain, while degree soundness is obtained by checking masked low-degree
relations at verifier-chosen indices.

\begin{algorithm}[H]
\caption{\textsf{DECS.DeriveGamma}(\(\mathsf{root},\eta,r,q\)) \(\to \Gamma\) (verifier and prover)}
\label{alg:decs-derivegamma}
\begin{algorithmic}[1]
\State Initialize counter \(ctr\gets 0\); let \(L\gets \big\lfloor \frac{2^{64}}{q}\big\rfloor\cdot q\).
\For{\(k=0\) to \(\eta-1\)}
  \For{\(j=0\) to \(r-1\)}
    \Repeat
      \State \(x \gets \textsf{SHA256}(\mathsf{root}\ \|\ \textsf{LE64}(ctr))\) interpreted as a 64-bit integer
      \State \(ctr\gets ctr+1\)
    \Until{\(x < L\)}
    \State \(\Gamma_{k,j}\gets x \bmod q\)
  \EndFor
\EndFor
\State \Return \(\Gamma\)
\end{algorithmic}
\end{algorithm}

\begin{algorithm}[H]
\caption{\textsf{DECS.VerifyEval}(\(\mathsf{root},\Gamma,R,\mathsf{open}\)) (verifier)}
\label{alg:decs-verifyeval}
\begin{algorithmic}[1]
\Require Domain size \(N\), modulus \(q\), mask count \(\eta\), committed row count \(r\).
\For{each opened index \(e\)}
  \State Verify Merkle authentication for \(e\) against \(\mathsf{root}\); reject on failure.
  \For{each mask row \(k<\eta\)}
    \State Check \(R_k^\star[e] \stackrel{?}{=} M_k(e) + \sum_{j=0}^{r-1}\Gamma_{k,j}\,P_j(e)\pmod q\).
  \EndFor
\EndFor
\State \Return accept
\end{algorithmic}
\end{algorithm}

\subsection{LVCS (Linear Vector Commitments)}
LVCS lifts DECS into a commitment to packed witness rows that supports succinct
openings of \(\Fq\)-linear forms of those rows. SmallWood uses LVCS both as a
batching layer and as the core of its hash-based PCS: the verifier requests
openings of linear forms corresponding to replay residuals, and DECS ensures
that the opened values come from low-degree committed row polynomials.

\begin{algorithm}[H]
\caption{\textsf{LVCS.CommitInitWithParams}(\(\mathrm{ringQ},\{r_j\}_{j<r},\ell,\mathrm{params}\)) \(\to (\mathsf{root},\mathsf{ProverKey})\) (prover)}
\label{alg:lvcs-commitinit}
\begin{algorithmic}[1]
\Require Ring \(R_q\) of size \(N\), packed rows \(r_j\in\Fq^{n_{\mathrm{cols}}}\), tail length \(\ell>0\).
\State Sample \(\mathrm{mask}_j\getsunif \Fq^{\ell}\) for each row \(j\) (HVZK tails).
\State Interpolate each row to \(P_j(X)\gets \textsf{interpolateRow}(\mathrm{ringQ},r_j,\mathrm{mask}_j,n_{\mathrm{cols}},\ell)\).
\State Commit to \(\{P_j\}\) via DECS and obtain Merkle root \(\mathsf{root}\).
\State Derive \(\Gamma\gets \textsf{DECS.DeriveGamma}(\mathsf{root},\eta,r,q)\) and cache values needed for openings.
\State \Return \((\mathsf{root},\mathsf{ProverKey})\).
\end{algorithmic}
\end{algorithm}

\begin{algorithm}[H]
\caption{\textsf{LVCS.CommitStep1}(\(\mathsf{root}\)) \(\to \Gamma\) (verifier)}
\label{alg:lvcs-commitstep1}
\begin{algorithmic}[1]
\State Store \(\mathsf{root}\) and derive \(\Gamma\gets \textsf{DECS.DeriveGamma}(\mathsf{root},\eta,r,q)\).
\State \Return \(\Gamma\).
\end{algorithmic}
\end{algorithm}

\begin{algorithm}[H]
\caption{\textsf{LVCS.CommitStep2}(\(\{R_k\}_{k<\eta}\)) \(\to\) accept/reject (verifier)}
\label{alg:lvcs-commitstep2}
\begin{algorithmic}[1]
\For{each mask row \(R_k\)}
  \State Reject if \(\deg(R_k) > d\) (degree guard before openings).
\EndFor
\State \Return accept.
\end{algorithmic}
\end{algorithm}

\begin{algorithm}[H]
\caption{\textsf{LVCS.EvalInitMany}(\(\mathrm{ringQ},\mathsf{ProverKey},\{\mathrm{Coeffs}_k\}_{k<m}\)) \(\to \{\bar v_k\}_{k<m}\) (prover)}
\label{alg:lvcs-evalinit}
\begin{algorithmic}[1]
\Require Linear forms \(\mathrm{Coeffs}_k\in\Fq^{r}\) for \(k=0,\dots,m-1\).
\State Let \(\ell\gets |\mathrm{mask}_j|\).
\For{\(k=0\) to \(m-1\)}
  \State Initialize \(\bar v_k\gets 0\in\Fq^{\ell}\).
  \State Set \(\bar v_k \gets \sum_{j=0}^{r-1} \mathrm{Coeffs}_k[j]\cdot \mathrm{mask}_j \in \Fq^{\ell}\).
\EndFor
\State \Return \(\{\bar v_k\}_{k<m}\).
\end{algorithmic}
\end{algorithm}

\begin{algorithm}[H]
\caption{\textsf{LVCS.ChooseE}(\(\ell,n_{\mathrm{cols}}\)) \(\to E\) (verifier)}
\label{alg:lvcs-choosee}
\begin{algorithmic}[1]
\Require Ring size \(N\); packed domain size \(n_{\mathrm{cols}}=|\Omega|\); masked slab indices \([n_{\mathrm{cols}},n_{\mathrm{cols}}+\ell)\).
\State Sample \(E\subset \{n_{\mathrm{cols}}+\ell,\dots,N-1\}\) uniformly without replacement with \(|E|=\ell\).
\State \Return \(E\).
\end{algorithmic}
\end{algorithm}

\begin{algorithm}[H]
\caption{\textsf{LVCS.EvalFinish}(\(\mathsf{ProverKey},E\)) \(\to \mathsf{Opening}\) (prover)}
\label{alg:lvcs-evalfinish}
\begin{algorithmic}[1]
\State \(\mathsf{open}\gets \mathsf{ProverKey}.\mathrm{DecsProver}.\textsf{EvalOpen}(E)\).
\State \Return \(\mathsf{Opening}\{\mathrm{DECSOpen}=\mathsf{open}\}\).
\end{algorithmic}
\end{algorithm}

\begin{algorithm}[H]
\caption{\textsf{LVCS.EvalStep2}(\(\{\bar v_k\},E,\mathsf{open},C,v_{\mathrm{targets}}\)) (verifier)}
\label{alg:lvcs-evalstep2}
\begin{algorithmic}[1]
\Require Coeffs \(C\in\Fq^{m\times r}\), public prefixes \(v_{\mathrm{targets}}\in(\Fq^{n_{\mathrm{cols}}})^{m}\).
\State Split \(\mathsf{open}\) into mask openings for \([n_{\mathrm{cols}},n_{\mathrm{cols}}+\ell)\) and tail openings at \(E\).
\State Verify the DECS openings (degree soundness and Merkle authentication); reject on failure.
\For{each \(k<m\) and each masked index \(i<\ell\)}
  \State Check \(\sum_{j=0}^{r-1} C[k,j]\cdot \mathsf{maskOpen.Pvals}[i][j] \stackrel{?}{=} \bar v_k[i]\pmod q\).
\EndFor
\State Reconstruct each batched polynomial \(Q_k(X)\) from its \(\Omega\)-prefix \(v_{\mathrm{targets}}[k]\) and mask tail \(\bar v_k\).
\State Check tail consistency of \(Q_k\) against the opened linear combinations at points in \(E\).
\State \Return accept.
\end{algorithmic}
\end{algorithm}

\subsection{LVCS and PCS/PIOP}
LVCS provides the linear opening queries needed by the polynomial oracle model,
while DECS provides degree soundness for the committed rows. SmallWood combines
the two into a transparent hash-based PCS used by the PACS and PIOP layers.

\subsection{What DECS/LVCS/PCS bind in ARC-SPRUCE}
\label{app:smallwood-binds}
At a high level, the SmallWood stack binds \emph{all committed witness rows of
one statement} under a single PCS root, and it enforces that every opened value
is consistent with one low-degree polynomial view of those rows. In our ARC
instantiation this binding matters in three ways:
\begin{itemize}[leftmargin=1.25em]
  \item \textbf{Row binding and low degree (DECS).}
  DECS ensures that every committed row is the low-degree polynomial matching
  the packed values on \(\Omega\) and the random tails on \(\Omega'\).
  \item \textbf{Replay binding across subsystems (LVCS).}
  The verifier reuses the same opened row values to replay every active
  constraint family. In issuance this covers carrier membership, decode
  bridges, the aligned commitment, the inverse-witness check, and the direct
  public-target identity \(T=B_0+U\vecc+Z\). In showing it covers
  signature shortness, the signature-side bridges to the full replay-image
  family \(\widehat T_0,\ldots,\widehat T_{B_{\mathrm{rep}}-1}\), the
  non-sign bridges to
  \(\widehat m_0,\ldots,\widehat m_{B_{\mathrm{rep}}-1}\),
  \(\widehat R_{0,0},\ldots,\widehat R_{0,B_{\mathrm{rep}}-1}\), and
  \(\widehat R_{1,0},\ldots,\widehat R_{1,B_{\mathrm{rep}}-1}\), the blockwise
  replay residuals, and the PRF companion constraints.
  \item \textbf{Single committed witness.}
  Because all committed rows live under one PCS root, the prover cannot use one
  hidden message for the hash relation and another for the PRF or the
  commitment. All replay checks refer to one coherent witness assignment.
\end{itemize}

\subsection{Openings requested in issuance and showing}
\label{app:smallwood-openings}
Both proofs follow the same pattern: the prover commits to all statement rows
as explicit-domain polynomials, the verifier derives Fiat--Shamir challenges,
and the prover opens the committed rows at verifier-chosen points. The verifier
then recomputes the active residuals from those openings and checks the
parallel replay identities at \(E'\) together with the aggregated
\(\Omega\)-sum identities.

\paragraph{Pre-sign witness surface and replay.}
The committed pre-sign rows decode to the source rows
\[
M,\ K,\ S,\ E,\ Z.
\]
The verifier replays the carrier membership relations, decode bridges, the aligned commitment equation
\[
\vecc = \mA_m M + \mA_k K + \mA_s S + E,
\]
the inverse-witness equation
\[
(B_3-R_1)\Had Z = 1,
\]
and the direct aligned target identity
\[
T = B_0 + U\vecc + Z.
\]
This appendix therefore records the compiled relation \(R_{\mathrm{Issue}}^{\mathrm{cmp}}\) only.

\paragraph{Post-sign witness surface and replay.}
The committed post-sign rows consist of:
\begin{enumerate}[leftmargin=1.25em]
    \item the signature limb rows;
    \item the carrier rows \(C^M\) and \(C^{\mathrm{ctr}}\);
    \item the decoded non-sign source rows \(M,K,R_0,R_1\);
    \item the full replay-image block families
    \[
    \widehat T_0,\ldots,\widehat T_{B_{\mathrm{rep}}-1},\qquad
    \widehat m_0,\ldots,\widehat m_{B_{\mathrm{rep}}-1},
    \]
    \[
    \widehat R_{0,0},\ldots,\widehat R_{0,B_{\mathrm{rep}}-1},\qquad
    \widehat R_{1,0},\ldots,\widehat R_{1,B_{\mathrm{rep}}-1};
    \]
    \item the PRF companion rows.
\end{enumerate}
The verifier replays:
\begin{enumerate}[leftmargin=1.25em]
    \item the shortness families on the cert-side signature surface;
    \item the carrier membership and decode families;
    \item the signature-side bridges
    \[
    \widehat T_b = \mathcal F_{\mathrm{rep},b}(\mA u)
    \]
    for every block \(b\);
    \item the non-sign transform bridges
    \[
    \widehat m_b = \mathcal F_{\mathrm{rep},b}(M+K),\qquad
    \widehat R_{0,b} = \mathcal F_{\mathrm{rep},b}(R_0),\qquad
    \widehat R_{1,b} = \mathcal F_{\mathrm{rep},b}(R_1)
    \]
    for every block \(b\);
    \item the blockwise replay residual
    \[
    (B_{3,b}^\Theta-\widehat R_{1,b})\Had\Bigl(\widehat T_b-B_{0,b}^\Theta-B_{1,b}^\Theta\Had \widehat m_b-B_{2,b}^\Theta\Had \widehat R_{0,b}\Bigr)
    = \mathcal F_{\mathrm{rep},b}(1)
    \]
    for every block \(b\); and
    \item the PRF checkpoint constraints.
\end{enumerate}
The theorem stated in Section~\ref{subsec:statement-strength} then transports
this full replay-image statement to the coefficient-domain BB-tran identity.

\paragraph{Replay-transform realization.}
The replay transform
\[
\mathcal F_{\mathrm{rep}} : \Rq \to \mathcal T
\]
is the exact public transform used by the replay basis, partitioned into the
replay blocks
\[
\mathcal F_{\mathrm{rep},0},\ldots,\mathcal F_{\mathrm{rep},B_{\mathrm{rep}}-1}.
\]
The bridge families above certify that the replay-image rows are exactly the
block projections of \(\mathcal F_{\mathrm{rep}}(\mA u)\),
\(\mathcal F_{\mathrm{rep}}(M+K)\), \(\mathcal F_{\mathrm{rep}}(R_0)\), and
\(\mathcal F_{\mathrm{rep}}(R_1)\). These are precisely the identities used in
Lemma~\ref{lem:showing-to-cleared}.

\paragraph{Proof of Lemma~\ref{lem:replay-transform}.}
The replay basis is chosen so that
\[
\mathcal F_{\mathrm{rep}} : \Rq \longrightarrow \mathcal T
\]
is the corresponding exact CRT/NTT transform of \(\Rq\). Its additivity and multiplicativity are the defining algebraic
properties of that transform, and injectivity follows because it is an isomorphism onto its image.

\paragraph{Proof of Lemma~\ref{lem:showing-to-cleared}.}
For every replay block \(b\), the source-to-replay bridges and the public lift definition give
\[
\widehat T_b = \mathcal F_{\mathrm{rep},b}(T),
\qquad
\widehat m_b = \mathcal F_{\mathrm{rep},b}(M+K),
\]
\[
\widehat R_{0,b} = \mathcal F_{\mathrm{rep},b}(R_0),
\qquad
\widehat R_{1,b} = \mathcal F_{\mathrm{rep},b}(R_1),
\]
\[
B_{i,b}^\Theta = \mathcal F_{\mathrm{rep},b}(B_i)
\qquad\text{for }i\in\{0,1,2,3\}.
\]
Substituting these identities into the blockwise replay residual yields
\[
\mathcal F_{\mathrm{rep},b}(B_3-R_1)
\Had
\Bigl(\mathcal F_{\mathrm{rep},b}(T)-\mathcal F_{\mathrm{rep},b}(B_0)-\mathcal F_{\mathrm{rep},b}(B_1(M+K))-
\mathcal F_{\mathrm{rep},b}(B_2R_0)\Bigr)
=
\mathcal F_{\mathrm{rep},b}(1)
\]
for every \(b\). By additivity and multiplicativity of \(\mathcal F_{\mathrm{rep},b}\), this is
\[
\mathcal F_{\mathrm{rep},b}\Bigl((B_3-R_1)\Had\bigl(T-B_0-B_1(M+K)-B_2R_0\bigr)\Bigr)
=
\mathcal F_{\mathrm{rep},b}(1)
\]
for every \(b\). Concatenating the replay blocks gives
\[
\mathcal F_{\mathrm{rep}}\Bigl((B_3-R_1)\Had\bigl(T-B_0-B_1(M+K)-B_2R_0\bigr)\Bigr)
=
\mathcal F_{\mathrm{rep}}(1).
\]
By Lemma~\ref{lem:replay-transform}, \(\mathcal F_{\mathrm{rep}}\) is injective. Hence
\[
(B_3-R_1)\Had\bigl(T-B_0-B_1(M+K)-B_2R_0\bigr)=1
\]
in \(\Rq\). Since the compiled showing relation also enforces \(T=\mA u\), this immediately implies
\[
(B_3-R_1)\Had\bigl(\mA u-B_0-B_1(M+K)-B_2R_0\bigr)=1
\]
in \(\Rq\).

\paragraph{The \(\Omega\)-sum check.}
The \(\Omega\)-sum check is active in both proofs. In issuance it aggregates the
carrier/decode checks together with the aligned commitment and public-target
identity. In showing it enforces the full signature-side bridge family together
with the full non-sign transform bridges. Thus the \(\Sigma_\Omega\) check is
part of the statement layer in both issuance and showing.

\subsection{PACS/PIOP interface}
SmallWood's PACS/PIOP layer checks that constraint polynomials vanish on
\(\Omega\) and at random points outside \(\Omega\), by committing to witness
rows and mask polynomials and then opening them at verifier-chosen
coordinates. The following algorithms summarize the high-level structure used in
our proofs.

\begin{algorithm}[H]
\caption{\textsf{PACS.PIOPCommitOracle}\(_{\mathrm{Prover}}(\textsf{params})\to \textsf{pcsCom}\)}
\label{alg:pacs-commitoracle}
\begin{algorithmic}[1]
\State Interpolate witness rows into \(P_1,\ldots,P_n\in\Fq[X]_{\le s+\ell'-1}\) with \(\ell'\) random blind evaluations outside \(\Omega\).
\State Sample mask polynomials \(M_1,\ldots,M_\rho\in\Fq[X]_{\le d_Q}\) such that \(\sum_{\omega\in\Omega} M_i(\omega)=0\) for all \(i\).
\State Commit to \(P\Vert M\) via PCS (implemented through LVCS/DECS); obtain \(\textsf{pcsCom}\).
\State \Return \(\textsf{pcsCom}\).
\end{algorithmic}
\end{algorithm}

\begin{algorithm}[H]
\caption{\textsf{PACS.PIOPBatchCombine}\(_{\mathrm{Prover}}(\textsf{pcsCom},\Gamma')\to Q\)}
\label{alg:pacs-batchcombine}
\begin{algorithmic}[1]
\State Using \(P,M\), form batched polynomials
\[
Q_i(X) \;=\; M_i(X) + \sum_j \Gamma'_{i,j}(X)\,F_j(X) + \sum_u \gamma'_{i,u}\,F'_u(X),
\]
where \(F_j\) and \(F'_u\) are the parallel and aggregated constraint polynomials evaluated on \(P\).
\State \Return \(Q=(Q_1,\ldots,Q_\rho)\).
\end{algorithmic}
\end{algorithm}

\begin{algorithm}[H]
\caption{\textsf{PACS.PIOPOpenAndCheck}(\(\textsf{pcsCom},Q\)) (verifier)}
\label{alg:pacs-openandcheck}
\begin{algorithmic}[1]
\State Sample a query set \(E'\subset S\) with \(|E'|=\ell'\) (Fiat--Shamir in NIZK).
\State Verify PCS openings of \(P\) and \(M\) at \(E'\) (via LVCS/DECS); reject on failure.
\State Check that each \(Q_i(e)\) matches its definition from \(M_i(e)\) and the constraint evaluations at \(e\in E'\).
\State Check the masked \(\Omega\)-sum test \(\sum_{\omega\in\Omega} Q_i(\omega)=0\) for all \(i\).
\State \Return accept/reject.
\end{algorithmic}
\end{algorithm}

In explicit-domain mode the prover also supplies the formal coefficient
descriptions of the active parallel and aggregated families. This ensures that
decode compositions, theta-lifted replay rows, and aggregated bridge
polynomials are replayed as polynomial identities at \(E'\).

\subsection{Fiat--Shamir compilation}
SmallWood compiles the interactive PACS/PIOP protocol to a NIZK using
Fiat--Shamir with grinding, following \cite{EPRINT:FenRiv25b}. In our
instantiation both issuance and showing use non-empty parallel and aggregated
families, so the verifier derives one batched challenge schedule and checks
both the replay identities at \(E'\) and the active \(\Omega\)-sum constraints.

\begin{algorithm}[H]
\caption{\textsf{PACS.ToNIZK} (Fiat--Shamir wrapper with grinding)}
\label{alg:pacs-tonizk}
\begin{algorithmic}[1]
\State Sample \(\textsf{salt}\getsunif\{0,1\}^{2\lambda}\).
\State Derive challenges by hashing the transcript with grinding counters to obtain required prefix properties.
\State Output PCS/LVCS/DECS openings and any required batched values; the verifier recomputes challenges and checks.
\end{algorithmic}
\end{algorithm}
